\subsection{고정 목표 NPC 운영 (NPC-ENG1/NPC-SCR1)}
Mira의 고정 목표는 항구 신호 복구 후 허가된 조석 단서를 제공하는 것이다. Python은 typed intent와 공개 관측으로 bounded BFS 계획을 계산한다. 획득·설치는 플레이어가, 허가된 \texttt{reveal\_hint}는 NPC가 수행한다. Godot은 session/epoch/request/goal/policy/revision/hash를 동기 검사한 뒤 기존 writer를 await 없이 호출한다. 결과와 새 관측이 다음 계획을 결정하므로 단일 후보의 validator-error repair와 구별된다.

NPC-ENG1은 실제 엔진 snapshot에서 시작한 R arm의 결정론적 engineering 사례 16개(assigned\_failure=2, completed=14)다. NPC-SCR1은 P/R/O 각 24개, 전체 72개 배정과 실패를 모두 보존한다(assigned\_failure=12, completed=60). P는 인증된 Codex CLI 실행 시도 24회, 반환 decision 24개, 배정 실패 4개를 기록했다. 이는 원격 과금 호출 횟수가 아니다. 요청 모델: gpt-5.6-sol; decision에 기록된 모델 식별자: gpt-5.6-sol. 후자는 설정된 selector이며 provider revision이 아니다. provider revision은 사용할 수 없고 null 식별자는 미상으로 남긴다. 전역 설정은 모델 식별 근거가 아니다. R은 작성된 규칙을 사용한다. O는 \texttt{same\_model\_family\_auxiliary}로 선언된 protocol auxiliary일 뿐이며 실제 모델 계열 일치는 unverified(검증되지 않음)이다. 독립 원본은 0개이고 O는 oracle이나 semantic gold가 아니다. 표~\ref{tab:npc-screening}과 그림~\ref{fig:npc-screening}은 bounded protocol-conformance 고정 사례의 기술적 계수이지 회복 효능·우월성·일반화 추정이 아니다. efficacy는 pending이며 토큰·비용은 미측정이다. 읽기 전용 CLI와 격리 작업 디렉터리를 사용했지만 tool-free 추론이나 적대적 기밀성은 검증하지 않았다.
검증된 사전 작성 stale-control ACK 거부는 arm별 P=4, R=4, O=4건(SCR-T01 (stale\_revision), SCR-T02 (stale\_policy))이며, 12개 인스턴스 모두 대응하는 음성 대조 gate가 모두 통과했다. 따라서 이 행들은 assigned failure로 분모에 남기되 method-performance 결과로 해석하지 않으며, 기타 assigned failure 0개는 미분류로 남긴다.
